\documentclass[10pt,a4paper]{amsart}
\usepackage[margin=19mm]{geometry}
\usepackage{amsmath,amssymb,mathtools}
\usepackage[scr=rsfs,cal=euler]{mathalfa}
\usepackage{xfrac}
\usepackage[unicode,hidelinks,bookmarks=false]{hyperref}
\newcommand{\ie}{i.e.\ }
\newcommand{\cf}{cf.\ }
\newcommand{\resp}{respectively\ }
\newcommand{\Subsection}{\subsection}
\providecommand{\nobreakdash}{\nobreak\hbox{-}}
\setlength{\emergencystretch}{2em}
\multlinegap0pt
\usepackage{bm}
\usepackage{bbm}
\usepackage{dsfont}
\usepackage{xspace}
\usepackage{cancel}
\usepackage{mathtools}
\usepackage{amsthm}
\makeatletter
\@ifundefined{definition}{\newtheorem{definition}{Definition}}{}
\@ifundefined{proposition}{\newtheorem{proposition}[definition]{Proposition}}{}
\makeatother
\newcommand{\Diracd}{\mathfrak{d}}
\newcommand{\R}{\mathbb{R}}
\newcommand{\ind}{1\!\!\mathrm{I}}
\newcommand{\rest}{\llcorner}
\renewcommand{\i}{\imath}
\title[Existence of spiral strategies for blocking fire spreading]{Existence of spiral strategies for blocking fire spreading}
\author{Stefano Bianchini, Martina Zizza}
\date{Source: arXiv:2508.05324v1 (CC BY 4.0)}
\begin{document}
\maketitle

\noindent\emph{Source and licence.} Existence of spiral strategies for blocking fire spreading. Version arXiv:2508.05324v1 (CC BY 4.0), distributed under the Creative Commons Attribution 4.0 International licence, as recorded in the arXiv licence metadata for this exact version and shown on the abstract page https://arxiv.org/abs/2508.05324v1. Full rights notice: licenses/arxiv-2508.05324-CC-BY-4.0.txt.\par\smallskip

\noindent\textit{Editorial scope.} Counted unit: the Proposition whose source label is \texttt{Prop:equa\_delta\_tilde\_r\_segm} in the article (source0.tex lines 4723--4739, source label \texttt{Prop:equa\_delta\_tilde\_r\_segm}), delivered together with the article's own complete original proof of it (source0.tex lines 4741--4807, one \texttt{proof} environment of 67 lines). No mathematics is written, repaired or re-derived: statement and proof are the article's own text line for line. Required context retained uncounted: Equa:delta\_theta\_segm\_1 (lines 4679--4682); Equa:delta\_P\_2P\_1\_segm (lines 4689--4696); Equa:segm\_equ\_tilde\_r (lines 4703--4710). Every definition, display and label of those lines is retained unchanged. Omitted from this package and not redistributed: 1--4678, 4683--4688, 4697--4702, 4711--4722, 4740--4740, 4808--12466. Nothing inside a retained excerpt is omitted or altered: every retained body is delivered exactly as the article's lines are written, so no omission receipt of any kind accompanies this package. Citations: the retained excerpts carry no citation command at all, so no bibliography entry of the article is transcribed and no reference is redistributed. Reference adaptation: none was needed and none was made; every reference inside the delivered unit resolves inside it, so no unresolved reference or citation is produced. Printed slips kept literal: no printed word, symbol, hypothesis or number of the counted statement or of its proof was altered. Layout and class: the article's own class is \texttt{amsart} with packages including array, color, float, esint, xcolor, ifpdf, graphicx, hyperref, graphics, siunitx, glossaries, glossary-mcols, aurical, amsfonts; the wrapper is the lane's pinned \texttt{amsart} editorial wrapper at 10pt on A4, which restates the theorem-like environments the retained text needs and adds the article's own macro definitions that the retained text needs (\texttt{\textbackslash Diracd}, \texttt{\textbackslash R}, \texttt{\textbackslash i}, \texttt{\textbackslash ind}, \texttt{\textbackslash rest}), with extra wrapper packages (bm, bbm, dsfont, xspace, cancel, mathtools). The delivered numbering is the unit's own. Assets: no retained span contains a figure environment, a graphics-inclusion command, a PDF-page inclusion, a table or a TikZ picture; no image file is copied into this package, the package declares no asset dependency and no image file is redistributed with it. Licence: version arXiv:2508.05324v1 of this article is distributed under the Creative Commons Attribution 4.0 International licence, as recorded in the arXiv licence metadata for this exact version and shown on the abstract page https://arxiv.org/abs/2508.05324v1. A full rights notice is delivered at licenses/arxiv-2508.05324-CC-BY-4.0.txt. Characters: every retained excerpt is pure ASCII, so the delivered engine, XeLaTeX with its default Latin Modern text font, reports no missing character.
\begin{equation}
\label{Equa:delta_theta_segm_1}
\frac{|P_2 - P_1| \delta \theta}{\delta s_0} = \frac{\cos(\bar \alpha - \theta) - \cos(\bar \alpha)}{\sin(\bar \alpha)}.
\end{equation}

\begin{equation}
\label{Equa:delta_P_2P_1_segm}
\begin{split}
\frac{\delta |P_2 - P_1|}{\delta s_0} - \frac{r(\phi_2)}{\sin(\bar \alpha)} \frac{\delta \theta}{\delta s_0} &= \cot(\bar \alpha) \frac{\cos(\bar \alpha) - \cos(\bar \alpha - \theta)}{\sin(\bar \alpha)} + \frac{\sin(\theta - \bar \alpha)}{\sin(\bar \alpha)} \\
&= \frac{\cos(\bar \alpha)^2 - \cos(\theta)}{\sin(\bar \alpha)^2} \\
&= - 1 + \frac{1 - \cos(\theta)}{\sin(\bar \alpha)^2}.
\end{split}
\end{equation}

\begin{equation}
\label{Equa:segm_equ_tilde_r}
\frac{d}{d\phi} \tilde r(\phi;s_0) = \cot(\bar \alpha) \tilde r(\phi;s_0) - \begin{cases}
R(\phi) & \phi_0 + \bar \theta - \bar \alpha \leq \phi < \phi_0 + 2\pi + \beta^-(\phi_0), \\
0 & \phi_0 + 2\pi + \beta^-(\phi_0) \leq \phi < \phi_0 + 2\pi + \bar \theta, \\
|P_2 - P_1| \Diracd_{\phi_0 + 2\pi + \bar \theta} + \frac{\tilde r(\phi - 2\pi - \bar \alpha)}{\sin(\bar \alpha)} & \phi \geq \phi_0 + 2\pi + \bar\theta.
\end{cases}
\end{equation}

\begin{proposition}
\label{Prop:equa_delta_tilde_r_segm}
For $\phi > \phi_2 = \phi_0 + \bar \theta - \bar \alpha$, the derivative 
$$
\delta \tilde r(\phi;s_0) = \lim_{\delta s_0 \searrow 0} \frac{\tilde r(\phi;s_0 + \delta s_0) - \tilde r(\phi;s_0)}{\delta s_0}
$$
satisfies the RDE on $\R$
\begin{equation*}
\frac{d}{d\phi} \delta \tilde r(\phi;s_0) = \cot(\bar \alpha) \delta \tilde r(\phi;s_0) - \frac{\delta \tilde r(\phi - 2\pi - \bar \alpha)}{\sin(\bar \alpha)} + S(\phi;s_0),
\end{equation*}
with source
\begin{align*}
S(\phi;s_0) &= \cot(\bar \alpha) \frac{1 - \cos(\theta)}{\sin(\bar \alpha)} \Diracd_{\phi_0 + \bar \theta - \bar \alpha} - \Diracd_{\phi_0 + 2\pi + \theta_0} \\
& \quad + \frac{\cos(\theta) - \cos(\bar \alpha)^2}{\sin(\bar \alpha)^2} \Diracd_{\phi_0 + 2\pi + \bar \theta} + \frac{\cos(\bar \alpha - \theta) - \cos(\bar \alpha)}{\sin(\bar \alpha)} \Diracd'_{\phi_0 + 2\pi + \bar \theta},
\end{align*}
where $\Diracd'$ is the distributional derivative of the Dirac's delta $\Diracd$.
\end{proposition}

\begin{proof}
%For shortness denote by $\phi_2$ the angle $\phi_0 + \bar \theta - \bar \alpha$.
%
We compute first the variation of the initial data of \eqref{Equa:segm_equ_tilde_r}: observing that the spiral $\tilde r(\phi;s_0)$ is saturated at $\phi_2$ and hence its tangent vector is $e^{\i (\phi_0 + \bar \theta)}$, one has up to second order terms for $\phi$ close to $\phi_2$ 
\begin{align*}
\tilde r(\phi;s_0 + \delta s_0) - \tilde r(\phi;s_0) &= \frac{e^{\i (\phi_0 + \bar \theta - \frac{\pi}{2})} \cdot \delta P_2}{\sin(\bar \alpha)} \\
&= \delta s_0 \frac{\sin(\theta)}{\sin(\bar \alpha)} - \frac{|P_2 - P_1| \delta \theta}{\sin(\bar \alpha)} \\
&= \delta s_0 \cot(\bar \alpha) \frac{1 - \cos(\theta)}{\sin(\bar \alpha)},
\end{align*}
which gives the first jump at $\phi_0 + \bar \theta - \bar \alpha$.

Being $R(\phi)$ the same up to $\phi_0 + 2\pi + \theta_0$, the RDE for $\delta \tilde r$ is actually the ODE
\begin{equation*}
\frac{d}{d\phi} \delta \tilde r(\phi) = \cot(\bar \alpha) \delta \tilde r(\phi)
\end{equation*}
for $\phi_0 + \bar \theta - \bar\alpha \leq \phi < \phi_0 + 2\pi + \theta_0$.

Using $R = Ds^-$, the variation at $\phi_0 + 2\pi + \theta_0$ is
\begin{align*}
\tilde r(\phi;s_0 + \delta s_0) - \tilde r(\phi;s_0) &= \cot(\bar \alpha) \big( \tilde r(\phi;s_0 + \delta s_0) - \tilde r(\phi;s_0) \big) - \big( s^-(\phi) - s^-(\phi_0 + 2\pi + \theta_0) \big) \\
&= \cot(\bar \alpha) \big( \tilde r(\phi;s_0 + \delta s_0) - \tilde r(\phi;s_0) \big) - \delta s_0.
\end{align*}
Hence $\delta \tilde r(\phi;s_0)$ has a jump of size $-1$ at the angle $\phi_0 + 2\pi + \theta_0$. %, due to the change of slope at the corner $P_1$.
%
From \eqref{Equa:segm_equ_tilde_r} it follows that
\begin{equation*}
\frac{d}{d\phi} \delta \tilde r(\phi) = \cot(\bar \alpha) \delta \tilde r(\phi)
\end{equation*}
for $\phi_0 + 2\pi + \theta_0 < \phi < \phi_0 + 2\pi + \bar \theta$.

At the angle $\phi_0 + 2\pi + \bar \theta$, computing the first order approximations for $0 \leq \phi - \phi_0 + 2\pi + \bar \theta \ll 1$ one gets up to second order terms
\begin{align*}
\tilde r(\phi;s_0 + \delta s_0) - \tilde r(\phi;s_0) &= \bigg[ \tilde r(\phi_0 + 2\pi + \bar\theta; s_0 + \delta s_0) + \cot(\bar \alpha) \tilde r(\phi_0 + 2\pi + \bar\theta;s_0 + \delta s_0) (\phi - \phi_0 - 2\pi - \bar \theta) \\
& \quad \quad - |P_2' - P_1'| \big( 1 + \cot(\bar \alpha) (\phi - \phi_0 - 2\pi - \bar \theta - \delta \bar \theta) \big) \ind_{\phi \geq \phi_0 + 2\pi + \bar \theta + \delta \bar \theta} \\
& \quad \quad - \frac{\tilde r(\phi_0 + \bar \theta - \bar \alpha;s_0 + \delta s_0)}{\sin(\bar \alpha)} (\phi - \phi_0 - 2\pi - \bar \theta - \delta \bar \theta) \bigg] \\
& \quad - \bigg[ \tilde r(\phi_0 + 2\pi + \theta-;s_0) + \cot(\bar \alpha) \tilde r(\phi_0 + 2\pi +\theta;s_0) (\phi - \phi_0 - 2\pi - \bar \theta) \\
& \quad \quad - |P_2 - P_1| \big( 1 + \cot(\bar \alpha) (\phi - \phi_0 - 2\pi - \bar \theta) \big) \ind_{\phi \geq \phi_0 + 2\pi + \bar \theta}  \\
& \quad \quad - \frac{\tilde r(\phi_0 + \bar \theta - \bar \alpha;s_0)}{\sin(\bar \alpha)} (\phi - \phi_0 - 2\pi - \bar \theta) \bigg] \\
&= \big( \tilde r(\phi_0 + 2\pi + \bar\theta;s_0 + \delta s_0) - \tilde r(\phi_0 + 2\pi + \bar\theta;s_0) \big) \big( 1 + \cot(\bar \alpha) (\phi - \phi_0 - 2\pi - \bar \theta) \big) \\
& \quad - \frac{\tilde r(\phi_0 + \bar \theta - \bar \alpha;s_0 + \delta s_0) - \tilde r(\phi_0 + \bar \theta - \bar \alpha;s_0)}{\sin(\bar \alpha)} (\phi - \phi_0 - 2\pi - \bar \theta) \\
& \quad + \bigg( - \delta |P_2 - P_1| + \frac{\tilde r(\phi_0 + \bar \theta - \bar \alpha;s_0)}{\sin(\bar \alpha)} \bigg) \ind_{\phi \geq \phi_0 + 2\pi + \bar \theta} + |P_2 - P_1| \ind_{[\phi_0 + 2\pi + \bar \theta,\phi_0 + 2\pi + \bar \theta + \delta \theta)}.
\end{align*}
For the above formula, we have used that the jumps $|P_2 - P_1|$ are in different angles and that the curvature term $- \frac{\tilde r(\phi_0+\bar \theta -\bar \alpha;s_0)}{\sin(\bar \alpha)}$ is applied starting from different angles. \\
Passing to the limit $\delta s_0 \searrow 0$ we obtain
\begin{align*}
\delta \tilde r(\phi;s_0) &= \delta \tilde r(\phi_0 + 2\pi + \bar\theta;s_0) \big( 1 + \cot(\bar \alpha) (\phi - \phi_0 - 2\pi - \bar \theta) \big) - \frac{\delta \tilde r(\phi_0 + \bar \theta - \bar \alpha;s_0)}{\sin(\bar \alpha)} (\phi - \phi_0 - 2\pi - \bar \theta) \\
& \quad + \bigg( - \frac{\delta |P_2 - P_1|}{\delta s_0} + \frac{\tilde r(\phi_0 + \bar \theta - \bar \alpha;s_0)}{\sin(\bar \alpha)} \frac{\delta \theta}{\delta s_0} \bigg) \ind_{\phi \geq \phi_0 + 2\pi + \bar \theta} + |P_2 - P_1| \frac{\delta \theta}{\delta s_0} \Diracd_{\phi_0 + 2\pi + \bar \theta}.
\end{align*}
Hence its derivative w.r.t. $\phi$ satisfies
\begin{align*}
\frac{d}{d\phi} \delta \tilde r(\phi;s_0) &= \cot(\bar \alpha) \delta \tilde r(\phi;s_0) - \frac{\delta \tilde r(\phi - 2\pi - \bar \alpha;s_0)}{\sin(\bar \alpha)} \\
& \quad + \bigg( - \frac{\delta |P_2 - P_1|}{\delta s_0} + \frac{\tilde r(\phi_0 + \bar \theta - \bar \alpha;s_0)}{\sin(\bar \alpha)} \frac{\delta \theta}{\delta s_0}  \bigg) \Diracd_{\phi_0 + 2\pi + \bar\theta} + |P_2 - P_1| \frac{\delta \theta}{\delta s_0} \Diracd'_{\phi_0 + 2\pi + \bar \theta},
\end{align*}
where $\Diracd'$ is the derivative of the Dirac's delta:
\begin{equation*}
\langle \Diracd'_x, f \rangle = - \frac{df(x)}{dx}.
\end{equation*}
Using \eqref{Equa:delta_theta_segm_1} and \eqref{Equa:delta_P_2P_1_segm} we obtain
\begin{align*}
- \frac{\delta |P_2 - P_1|}{\delta s_0} + \frac{\bar r(\phi_0 + \bar \theta - \bar \alpha;s_0)}{\sin(\bar \alpha)} \frac{\delta \theta}{\delta s_0} &= \cot(\bar \alpha) \frac{\cos(\bar \alpha - \theta) - \cos(\bar \alpha)}{\sin(\bar \alpha)} - \frac{\sin(\theta - \bar \alpha)}{\sin(\bar \alpha)} \\
&= \frac{\cos(\theta) - \cos(\bar \alpha)^2}{\sin(\bar \alpha)^2},
\end{align*}
\begin{equation*}
|P_2 - P_1| \frac{\delta \theta}{\delta s_0} = \frac{\cos(\bar \alpha - \theta) - \cos(\bar \alpha)}{\sin(\bar \alpha)}.
\end{equation*}
It is now easy to verify that the derivative $\delta \tilde r(\phi;s_0) \rest_{\phi_0 + \bar \theta - \bar \alpha}$ satisfies the RDE in the statement, since in the remaining part of $\R$ it solves the linear RDE of the saturated spiral.
\end{proof}

\end{document}
